\documentclass[11pt,a4paper]{article}
\usepackage[margin=2.2cm]{geometry}
\usepackage[T1]{fontenc}
\usepackage[utf8]{inputenc}
\usepackage{lmodern}
\usepackage{booktabs}
\usepackage{longtable}
\usepackage{tabularx}
\usepackage{hyperref}
\usepackage{enumitem}
\usepackage{xcolor}
\usepackage{titlesec}
\usepackage{fancyhdr}
\usepackage{graphicx}
\usepackage{amsmath}

\hypersetup{colorlinks=true,linkcolor=black,urlcolor=blue!50!black}
\pagestyle{fancy}
\fancyhf{}
\fancyhead[L]{\small AI Museum of the Future}
\fancyhead[R]{\small Technical reference --- models, metrics, hyperparameters}
\fancyfoot[C]{\thepage}
\renewcommand{\headrulewidth}{0.4pt}

\title{\textbf{AI Museum of the Future}\\[0.4em]
\large Technical reference: models, hyperparameters, metrics\\
and end-to-end path to the museum exhibition}
\author{Team project --- ESPRIT / Deep Learning\\
\textit{Document for oral defence (professor)}}
\date{October 2026}

\begin{document}
\maketitle
\begin{abstract}
\noindent
This note lists every major model used (or planned) in the project, the
hyperparameters that control retrieval and agents, the measured metrics, and
the full pipeline from archival documents to a curated museum exhibition.
\textbf{Deployed today:} text LLM (Groq), local embeddings (MiniLM), vision
understanding (Qwen). \textbf{Image generation} is not in production; we
document \textbf{Flux.1} as the technical candidate for optional hall
illustrations.
\end{abstract}

\tableofcontents
\vspace{1em}

%% =====================================================================
\section{End-to-end path: from request to museum}
%% =====================================================================

The system turns a visitor (or curator) request into a multi-section exhibition
grounded in documentary evidence.

\begin{enumerate}[leftmargin=*,itemsep=2pt]
  \item \textbf{Input.} Country / year range / topic (default Tunisia, 2020--2026),
        via Flask dashboard, visitor UI, or n8n webhook.
  \item \textbf{Knowledge base.} Documents collected, chunked, embedded, stored in
        Supabase (Postgres + pgvector + full-text search).
  \item \textbf{Hybrid RAG.} For each specialist: vector search + keyword search,
        fused, filtered, diversified; optional LLM rerank (off by default).
  \item \textbf{Five specialists (parallel).} Historian, Sociologist, Technology,
        Culture~\&~Art, Visual --- each with its own search queries and prompt.
  \item \textbf{Multimodal (Visual).} Up to $N$ images are described by a vision
        model; descriptions become evidence lines \texttt{vision:\ldots}.
  \item \textbf{Governance.} Peer reviewer $\rightarrow$ Curator (4--6 sections)
        $\rightarrow$ Fact-checker (unsupported claims trigger one curator rewrite).
  \item \textbf{Persistence.} Run saved as \texttt{runs/<run\_id>/run.json}
        (prompts, evidence, insights, timings).
  \item \textbf{Experience layer.} Exposition page, visitor 3D halls (corpus
        \texttt{image\_url} on walls), optional n8n delivery / Binary image nodes.
\end{enumerate}

\begin{center}
\small
\texttt{collectors $\rightarrow$ chunking $\rightarrow$ MiniLM embeddings $\rightarrow$ Supabase}\\[2pt]
\texttt{$\downarrow$ hybrid retrieval}\\[2pt]
\texttt{5 specialists (+ vision for Visual) $\rightarrow$ review $\rightarrow$ curator $\rightarrow$ fact-check}\\[2pt]
\texttt{$\downarrow$}\\[2pt]
\texttt{Flask / visitor UI / n8n $\rightarrow$ museum exhibition}
\end{center}

%% =====================================================================
\section{Models catalogue}
%% =====================================================================

\subsection{Models deployed in the running system}

\begin{longtable}{@{}p{3.2cm}p{4.2cm}p{2.2cm}p{5.2cm}@{}}
\toprule
\textbf{Role} & \textbf{Model / component} & \textbf{Provider} & \textbf{Env / notes} \\
\midrule
\endhead
Text generation \& reasoning (all specialists + reviewer + curator + fact-checker by default) &
\texttt{openai/gpt-oss-120b} &
Groq (OpenAI-compatible API) &
\texttt{LLM\_MODEL}; \texttt{LLM\_BASE\_URL=https://api.groq.com/openai/v1} \\
\addlinespace
Optional stronger curator &
same or override &
Groq / other &
\texttt{CURATOR\_MODEL} \\
\addlinespace
Vision / multimodal \textbf{understanding} &
\texttt{qwen/qwen3.8-27b} &
Groq (vision-capable) &
\texttt{VISION\_MODEL}; describes real corpus images \\
\addlinespace
Optional retrieval rerank &
defaults to \texttt{LLM\_MODEL} &
same &
\texttt{RERANK\_MODEL}; off by default in agent searches \\
\addlinespace
Embeddings (local) &
\texttt{paraphrase-multilingual-MiniLM-L12-v2} &
FastEmbed (on-device) &
384 dimensions; $\sim$120\,MB download once; no API key \\
\addlinespace
Vector DB &
pgvector in Supabase &
Supabase &
\texttt{vector(384)}; RPC \texttt{match\_chunks} + FTS \\
\addlinespace
Orchestration &
n8n workflow &
self-hosted / cloud &
start $\rightarrow$ poll $\rightarrow$ summary / museum URLs \\
\addlinespace
Application &
Flask + static frontends &
Python &
visitor, dashboard, Agent lab, exposition API \\
\bottomrule
\end{longtable}

\subsection{Image generation candidate (not deployed): Flux.1}

\begin{itemize}[leftmargin=*]
  \item \textbf{Name:} Flux.1 (Black Forest Labs) --- modern text-to-image model.
  \item \textbf{Role in our architecture (planned / oral answer):} optional
        \emph{illustration} of museum halls or atmospheric panels when no corpus
        photo exists --- always labelled as generated, never mixed silently with
        documentary evidence.
  \item \textbf{Why Flux.1 is a good technical answer:} strong prompt adherence
        and image quality vs older Stable Diffusion baselines; available via API
        or GPU inference; industry-standard choice in 2025--2026 demos.
  \item \textbf{How Flux.1 works (high level):}
        \begin{enumerate}[nosep]
          \item Encode the text prompt (and optional style constraints).
          \item Sample from noise in a latent space.
          \item Iteratively transform / denoise (flow-matching / diffusion-family
                dynamics) toward an image matching the prompt.
          \item Decode the latent to a pixel image.
        \end{enumerate}
  \item \textbf{Why we do not use it in production today:} free-tier credits,
        documentary integrity (museum of memory prefers real \texttt{image\_url}),
        and multimodal priority = \emph{understanding} captions vs pixels.
\end{itemize}

\noindent\textit{One-line oral answer:}
``For generation we would use Flux.1; for understanding we already use Qwen vision;
today the walls show real archive images.''

%% =====================================================================
\section{What each agent does (same LLM, different prompts)}
%% =====================================================================

All specialists share \texttt{agents/base\_agent.py}: gather evidence $\rightarrow$
LLM call (JSON) $\rightarrow$ citation filter $\rightarrow$ insight or abstain.

\begin{longtable}{@{}p{2.6cm}p{3.4cm}p{8.5cm}@{}}
\toprule
\textbf{Agent} & \textbf{Model} & \textbf{Focus / special knobs} \\
\midrule
\endhead
Historian & \texttt{LLM\_MODEL} & Timelines, events, long memory; may write in French. \\
Sociologist & \texttt{LLM\_MODEL} & Society, lifestyles; own search + \textbf{HyDE} (hypothetical document expansion, cached under \texttt{HYDE\_CACHE\_DIR}). \\
Technology & \texttt{LLM\_MODEL} & Innovations, tools, implications; dedicated search queries \& prompt v5. \\
Culture \& Art & \texttt{LLM\_MODEL} & Arts, media, symbols; category-aware search. \\
Visual & \texttt{LLM\_MODEL} + \texttt{VISION\_MODEL} & Images first: describe up to \texttt{VISUAL\_MAX\_IMAGES} (default 3), then interpret. Search category \texttt{image}. \\
Peer reviewer & \texttt{LLM\_MODEL} & Agreements, tensions, weak findings across insights. \\
Curator & \texttt{LLM\_MODEL} or \texttt{CURATOR\_MODEL} & Title, concept, 4--6 sections, artifacts from specialists only; one retry with feedback on validation failure. \\
Fact-checker & \texttt{LLM\_MODEL} & Narrative vs cited evidence; unsupported $\rightarrow$ one curator rewrite. \\
\bottomrule
\end{longtable}

\noindent\textbf{Hard citation rule:} every finding must contain at least one
\texttt{[chunk\_id]} present in that agent's evidence; uncited findings are
dropped; if none remain, the agent abstains and appears in
\texttt{missing\_perspectives}.

%% =====================================================================
\section{Hyperparameters and configuration knobs}
%% =====================================================================

\subsection{LLM / runtime}

\begin{tabular}{@{}llp{7.5cm}@{}}
\toprule
\textbf{Parameter} & \textbf{Default} & \textbf{Effect} \\
\midrule
\texttt{LLM\_MODEL} & \texttt{openai/gpt-oss-120b} & Text reasoning / generation \\
\texttt{VISION\_MODEL} & \texttt{qwen/qwen3.8-27b} & Image understanding \\
\texttt{CURATOR\_MODEL} & (unset $=$ LLM) & Optional stronger curator \\
\texttt{LLM\_MAX\_CONCURRENCY} & 3 & Parallel LLM calls (TPM-friendly) \\
\texttt{LLM\_ATTEMPTS} & 4 & Retries with provider-aware wait \\
Response format & JSON object & Structured agent outputs \\
Specialist timeout & 120\,s (often 400\,s free tier) & Pipeline \texttt{--timeout} \\
Full run LLM calls & $\sim$15--20 & Quota planning on Groq free tier \\
\bottomrule
\end{tabular}

\subsection{Embeddings \& chunking}

\begin{tabular}{@{}llp{7.5cm}@{}}
\toprule
\textbf{Parameter} & \textbf{Default / choice} & \textbf{Rationale} \\
\midrule
Embedding model & MiniLM multilingual L12 & Multilingual FR/EN/AR corpus \\
Dimension & 384 & Matches \texttt{vector(384)} schema \\
Effective window & 128 tokens & Model ignores tokens beyond 128 \\
Chunking strategy & Sentence packing to 128 + title prefix & Chosen over fixed-512 (57\% text unseen) \\
Chunking MRR@10 (selection) & 0.47 & Passage-level offline comparison \\
\bottomrule
\end{tabular}

\subsection{Hybrid retrieval (\texttt{SearchConfig})}

\begin{tabular}{@{}llp{7.2cm}@{}}
\toprule
\textbf{Parameter} & \textbf{Default} & \textbf{Meaning} \\
\midrule
\texttt{candidate\_pool} & 40 & Candidates per channel before fuse/filter \\
\texttt{statistics\_weight} & 0.35 & Down-weight World Bank-style stats rows \\
\texttt{max\_statistics} & 1 & At most one stats row in a result set \\
\texttt{duplicate\_threshold} & 0.6 & Token Jaccard near-duplicate cut \\
\texttt{rerank\_pool} & 20 & Candidates sent to LLM reranker \\
\texttt{relax\_country} & true & Top-up outside country if too few hits \\
\texttt{rerank\_agent\_searches} & false & LLM rerank off for agent own-search \\
Default year filter & 2020--2026 & Exhibition default window \\
\texttt{keep\_undated} & true & Many docs lack \texttt{published\_at} \\
Result diversity & 1 chunk / document & Avoid same-doc flooding \\
\bottomrule
\end{tabular}

\subsection{Multimodal}

\begin{tabular}{@{}llp{7.5cm}@{}}
\toprule
\textbf{Parameter} & \textbf{Default} & \textbf{Meaning} \\
\midrule
\texttt{VISUAL\_MAX\_IMAGES} & 3 & Vision LLM calls per Visual run \\
Download User-Agent & required & Wikimedia returns 403 without it \\
Cache & \texttt{multimodal/cache/} & Avoid re-describing same URL \\
Output & structured JSON + text line & Fed as \texttt{vision:} evidence; optional DB \texttt{vision\_*} \\
\bottomrule
\end{tabular}

\subsection{Curator validation}

\begin{itemize}[nosep]
  \item Section count must be in $[4,6]$.
  \item Non-empty title / concept / narratives.
  \item Artifacts must come from specialists (no invented items).
  \item One retry with explicit feedback; then \texttt{CuratorValidationError}.
\end{itemize}

\subsection{Flux.1 (if enabled later) --- typical generation knobs}

\begin{tabular}{@{}llp{7.2cm}@{}}
\toprule
\textbf{Knob} & \textbf{Typical range} & \textbf{Role} \\
\midrule
Prompt (+ negative) & free text & Subject, style, documentary vs illustrative \\
Steps / guidance & model-dependent & Quality vs latency / cost \\
Resolution & e.g.\ 1024$\times$1024 & Hall panel size \\
Seed & fixed for demos & Reproducibility \\
\label & ``AI illustration'' & Ethics / museum honesty \\
\bottomrule
\end{tabular}

%% =====================================================================
\section{Metrics and evaluation results}
%% =====================================================================

\subsection{Retrieval (documented in \texttt{evaluation/})}

Test set: \textbf{16 queries}, $k=5$, term-rule relevance proxy
(\texttt{retrieval\_cases.json}).

\paragraph{Run 2026-10-14 (README table --- often cited in pitch).}
\begin{center}
\begin{tabular}{@{}lcccc@{}}
\toprule
Method & Precision@5 & Hit@5 & MRR & Stats share \\
\midrule
Vector only & 0.41 & 0.75 & 0.52 & 0.11 \\
Hybrid & 0.68 & 0.94 & 0.83 & 0.00 \\
Hybrid + LLM rerank & \textbf{0.72} & \textbf{1.00} & \textbf{0.92} & 0.00 \\
\bottomrule
\end{tabular}
\end{center}

\paragraph{Update 2026-10-15 (full reranks completed).}
Vector $0.50 / 0.88 / 0.61$; hybrid $0.70 / 1.00 / 0.94$; hybrid+rerank
$0.56 / 1.00 / 0.75$. With every rerank executed, reranking \textbf{hurt}
precision --- hence \textbf{rerank off by default}.

\paragraph{Metric definitions.}
\begin{itemize}[nosep]
  \item \textbf{Precision@$k$:} fraction of top-$k$ results that match relevance rules.
  \item \textbf{Hit@$k$:} fraction of queries with $\ge 1$ relevant result in top-$k$.
  \item \textbf{MRR:} mean of $1/\text{rank}$ of the first relevant result.
\end{itemize}

\noindent\textit{Caveat:} Tunisia-specific precision stays low ($\sim$0.2--0.4): few
Tunisia chunks ($\sim$75, many statistics). Better retrieval cannot invent missing archives.

\subsection{Chunking selection}

Offline comparison on long documents / passage questions (EN/FR/AR): sentence packing
to the embedder's 128-token window + title prefix selected (MRR@10 $\approx 0.47$
in D1 draft notes).

\subsection{Exhibition run metrics (\texttt{evaluate\_run})}

\begin{tabular}{@{}lp{10cm}@{}}
\toprule
\textbf{Metric} & \textbf{Meaning} \\
\midrule
coverage & Specialists with cited findings / 5 \\
abstained / errors & No valid citations / provider timeout \\
findings\_dropped\_uncited & Guardrail: uncited claims removed \\
citation\_survival & kept / (kept + dropped) \\
own\_search\_results\_per\_agent & Evidence each agent retrieved itself \\
review\_* & Peer-review agreements, tensions, weak findings \\
fact\_check\_supported & Sections fully supported by cited text \\
faithfulness (\texttt{--judge}) & LLM judge: finding supported by its citations \\
llm\_calls / timings\_s & Cost and latency of the run \\
\bottomrule
\end{tabular}

\subsection{Multimodal qualitative metric}

Caption ``Protest in Mexico City'' vs vision: beach/road scene, no protest
$\rightarrow$ \texttt{caption\_mismatch}. Demonstrates why vision understanding
matters beyond text RAG.

%% =====================================================================
\section{Multimodal module --- image understanding (not generation)}
%% =====================================================================

\begin{enumerate}[leftmargin=*]
  \item Read \texttt{image\_url} from document metadata.
  \item Download image (browser User-Agent).
  \item Call \texttt{VISION\_MODEL} $\rightarrow$ structured JSON description.
  \item Disk-cache; attach \texttt{vision:} line to Visual evidence.
  \item Optionally write \texttt{vision\_*} into Supabase so hybrid search can
        match visual content, not only captions.
\end{enumerate}

\noindent Exposition UI can show ``what the AI sees'' via
\texttt{POST /api/exhibition/<run\_id>/vision}.

%% =====================================================================
\section{Orchestration and museum delivery}
%% =====================================================================

\begin{itemize}[leftmargin=*]
  \item \textbf{Flask API:} start / cancel / poll / list runs; vision describe;
        Agent lab \texttt{POST /api/agents/<name>/analyze}.
  \item \textbf{n8n:} \texttt{Exhibition\_pipeline\_end\_to\_end.json} ---
        webhook/schedule $\rightarrow$ start pipeline $\rightarrow$ poll
        ($\sim$60\,s) $\rightarrow$ summary + museum / halls URLs; Binary
        download of corpus images for walls.
  \item \textbf{Visitor UX:} warm public frontend + 3D rooms (specialist
        galleries) + immersive visit; login for protected routes
        (demo \texttt{visiteur} / \texttt{musee2126}).
  \item \textbf{Dashboard:} Agent lab (test each specialist), exhibition board,
        data browser.
\end{itemize}

%% =====================================================================
\section{Trust, security, and limits}
%% =====================================================================

\begin{itemize}[leftmargin=*]
  \item Traceability via mandatory \texttt{[chunk\_id]} citations.
  \item Structured JSON + curator / fact-check validators.
  \item Secrets in \texttt{.env} (never committed); optional Bearer tokens for APIs.
  \item Free-tier TPM can block curator rewrite after fact-check.
  \item Citation $\neq$ proof: relevance and fact-check remain necessary.
  \item Image generation (Flux.1) kept optional to protect documentary honesty.
\end{itemize}

%% =====================================================================
\section{Quick oral cheat-sheet}
%% =====================================================================

\begin{enumerate}[leftmargin=*]
  \item \textbf{Text brain:} Groq \texttt{gpt-oss-120b}.
  \item \textbf{Memory:} MiniLM 384-d + Supabase hybrid RAG.
  \item \textbf{Eyes:} Qwen vision --- understand photos.
  \item \textbf{Paintbrush (future):} Flux.1 --- generate illustrations only.
  \item \textbf{Best retrieval numbers to quote carefully:} hybrid+rerank
        0.72 / 1.00 / 0.92 (14 Oct) \emph{or} hybrid 0.70 / 1.00 / 0.94 with
        rerank off (15 Oct finding).
  \item \textbf{Museum result:} 4--6 curated sections + corpus images on walls,
        not AI-invented history.
\end{enumerate}

\vspace{1.5em}
\noindent\rule{\textwidth}{0.4pt}\\
{\small\textit{Sources in repository:} \texttt{README.md}, \texttt{agents/llm\_config.py},
\texttt{rag/vector\_store/hybrid\_retriever.py}, \texttt{multimodal/README.md},
\texttt{evaluation/README.md}, \texttt{docs/REPORT\_D1.md}. Confirm live
\texttt{.env} model names before the defence if the team changed providers.}

\end{document}
